\documentclass[11pt,a4paper]{article}
\usepackage[utf8]{inputenc}
\usepackage[T1]{fontenc}
\usepackage[french]{babel}
\usepackage{amsmath,amssymb,amsthm}
\usepackage{graphicx}
\usepackage{booktabs}
\usepackage{hyperref}
\usepackage{geometry}
\geometry{margin=2.5cm}

\title{\textbf{SOCRATES Phase 2, Étape 2 :} \\
\large Calcul Poly-Algébrique --- Estimation de Dimension par Hypergraphes, \\
un Banc d'Essai de Dix Problèmes Physiques, et Comparaison aux Méthodes Classiques}
\author{Xavier Callens \\ Socrate AI Lab}
\date{\today}

\begin{document}
\maketitle

\begin{abstract}
Ce rapport documente la Phase 2 (Étape 2) du programme Socrate AI Lab : la
construction et l'évaluation rigoureuse d'une boîte à outils computationnelle
de « Calcul Poly-Algébrique » pour la physique fondée sur les hypergraphes,
suivant la proposition de Wolfram selon laquelle l'espace, le temps et la
physique connue pourraient émerger de la réécriture discrète d'un
hypergraphe pré-géométrique. Parmi quatre concepts de formalisation
envisagés, deux ont été retenus pour implémentation sur des critères de
faisabilité~: un calcul de recherche de règles d'arité supérieure (le
\emph{Calcul Poly-Algébrique} proprement dit) et un estimateur de dimension
émergente (les \emph{Variables de Dimension Pré-Géométrique}). L'estimateur
de dimension a été appliqué à dix problèmes physiques connus dont la
dimension cible est vérifiable indépendamment --- résultats exacts pour des
orbites fermées et un théorème rigoureux, et valeurs de référence issues de
la littérature (Grassberger--Procaccia, 1983) pour deux attracteurs
chaotiques classiques. Le premier banc d'essai a produit un résultat
affiché de « 10/10 réussis » qu'un audit ultérieur a montré être largement
trompeur~: six des dix réussites provenaient d'un régime mathématiquement
dégénéré, et aucun des deux résultats chaotiques n'était convergé en nombre
de points. Deux défauts logiciels réels et indépendants ont été identifiés
au cours du processus et sont corrigés ici, avec tests de non-régression et
mesures avant/après, incluant une amélioration de performance d'environ
27$\times$. Une véritable méthode de référence classique (la méthode de
somme de corrélation de Grassberger--Procaccia) ainsi qu'un dispositif de
comparaison équitable ont ensuite été construits afin que l'affirmation
« meilleur que l'approche traditionnelle » devienne une assertion
vérifiable plutôt qu'une simple déclaration. Une nouvelle comparaison
systématique sur les dix problèmes, visant un taux de réussite majoritaire
face à cette référence, était en cours au moment de la rédaction et est
rapportée ici comme un état intermédiaire, non comme un résultat final ---
conformément à la règle permanente de ce programme selon laquelle une
mesure non convergée ou incomplète doit être signalée comme telle, jamais
lissée en une conclusion prématurée.
\end{abstract}

\tableofcontents

\section{Introduction et Motivation}

Le projet de physique de Stephen Wolfram propose que l'espace, le temps et
les lois connues de la physique ne sont pas fondamentaux, mais émergent de
l'application discrète et répétée de règles de réécriture locales sur un
hypergraphe pré-géométrique. Formaliser cette hypothèse de façon
computationnelle nécessite de nouveaux outils mathématiques et logiciels,
de manière analogue aux ruptures paradigmatiques historiques en physique
qui ont nécessité une nouvelle syntaxe mathématique (l'algèbre pour la
quantité inconnue, le calcul infinitésimal pour le changement instantané,
les nombres complexes pour les solutions « hors de la droite réelle »).

Quatre concepts ont été envisagés comme outils formels candidats pour ce
programme~:

\begin{enumerate}
\item \textbf{Calcul Poly-Algébrique} (algèbre d'arité supérieure) :
  étendre le « résoudre pour une quantité inconnue » de l'algèbre classique
  à « résoudre pour une relation ou une règle de réécriture inconnue
  d'arité supérieure ».
\item \textbf{Cohomologie Branchiale} : un invariant topologique du graphe
  multivoie (branchant) de tous les historiques de réécriture possibles,
  proposé comme l'analogue discret de l'interférence quantique.
\item \textbf{Opérateurs Ruliaux} : des dérivées discrètes mesurant le taux
  de déformation topologique par étape de réécriture, proposées comme
  analogue discret de la densité d'énergie/masse.
\item \textbf{Variables de Dimension Pré-Géométrique} : traiter la
  dimension de l'espace émergent comme une quantité fluide et localement
  variable, déduite de la façon dont le volume de boule de l'hypergraphe
  croît avec le rayon.
\end{enumerate}

\subsection{Justification de la sélection}

Les concepts \#1 et \#4 ont été retenus pour implémentation. Le concept
\#4 (estimation de dimension par croissance de volume) est une technique
bien établie déjà utilisée en théorie des ensembles causaux et dans les
travaux de Wolfram lui-même --- immédiatement implémentable sans
mathématiques nouvelles. Le concept \#1 (un calcul de recherche de règles)
est le substrat dont dépend implicitement tout autre concept~: on ne peut
mesurer une dimension, une dérivée rulialle ou un invariant branchial sans
disposer au préalable d'un langage rigoureux pour les hypergraphes, les
motifs et les réécritures. Les concepts \#2 et \#3 ont été jugés les moins
réalisables~: la \emph{Cohomologie Branchiale} en particulier n'a, à ce
jour, aucune définition mathématique opérationnelle au-delà d'un nom
évocateur --- aucun complexe de chaînes, aucun opérateur de bord n'a été
spécifié nulle part dans la littérature connue de ce programme, et son
implémentation aurait nécessité d'inventer les mathématiques ex nihilo
plutôt que d'assembler des outils existants, en tension avec la discipline
de ce programme qui consiste à ne jamais affirmer plus qu'un calcul ne le
prouve réellement.

\section{Méthodes}

\subsection{Représentation et réécriture d'hypergraphes}

Un hypergraphe est représenté comme un multi-ensemble immuable
d'hyperarêtes ordonnées sur un ensemble de nœuds implicite (les nœuds
n'existent qu'en vertu de leur participation à une relation, suivant la
convention du Modèle de Wolfram). Les règles de réécriture sont spécifiées
par un motif de partie gauche (un tuple d'arêtes-motifs sur des variables
de motif) et un remplacement de partie droite~; la correspondance est
effectuée par recherche avec retour arrière sur les affectations cohérentes
des variables de motif (filtrage de sous-hypergraphe). Deux modes
d'évolution sont implémentés~: un mode à historique unique qui applique
toutes les correspondances non chevauchantes à chaque génération (la règle
de mise à jour standard du Modèle de Wolfram), et un mode multivoie qui se
ramifie sur chaque mise à jour à correspondance unique individuelle,
produisant le graphe d'états branchant complet du système de réécriture
abstrait.

\subsection{Estimation de dimension émergente}

Pour un hypergraphe aplati en son graphe d'adjacence (chaque hyperarête
devient une clique sur ses nœuds constitutifs), la boule de rayon $r$
autour d'un nœud est l'ensemble des nœuds atteignables en $r$ sauts. Si le
graphe se comporte comme un réseau de dimension $d$ à une certaine échelle,
le volume de boule croît comme $|B(r)| \sim r^d$. L'estimateur ajuste la
\emph{taille de coquille} ($dV(r) = |B(r)| - |B(r-1)|$), et non le volume
cumulé, en espace log-log~: le volume cumulé sur un réseau est affine et
non homogène (un chemin unidimensionnel a $|B(r)| = 2r+1$ depuis un point
intérieur, et non $r^1$), ce qui biaise un ajustement naïf sur le volume
lui-même bien en dessous de la dimension réelle à tout rayon
computationnellement praticable. La taille de coquille n'a pas ce décalage
additif (la coquille d'un chemin est la constante exacte 2~; celle d'une
grille 2D est exactement $4r$), de sorte que l'estimateur corrigé retourne
\emph{exactement} 1,0 et 2,0 sur ces réseaux connus, non pas seulement
approximativement.

\subsection{Des trajectoires continues aux hypergraphes}

Une trajectoire physique continue (un ensemble de points dans l'espace des
phases) est convertie en hypergraphe via un graphe de proximité des $k$
plus proches voisins~: chaque point devient un nœud, et une arête relie
chaque point à ses $k$ plus proches voisins (symétrisé). La distance de
graphe sur un graphe de $k$-PPV suffisamment dense approxime la distance
géodésique sur la variété sous-jacente, la justification standard de cette
classe d'estimateur de dimension (Grassberger \& Procaccia, 1983~; la même
construction sous-tend l'apprentissage de variétés de type Isomap).

\subsection{La méthode de référence traditionnelle}

Afin que « meilleur que l'approche traditionnelle » devienne une
affirmation vérifiable, l'intégrale de corrélation classique de
Grassberger--Procaccia a été implémentée comme méthode de référence
explicite~:
\[
C(r) = \frac{2}{N(N-1)} \#\{(i,j) : i<j,\ |x_i - x_j| < r\}, \qquad
D_2 = \frac{d \log C(r)}{d \log r}
\]
Cette référence et l'estimateur par croissance de coquilles reposent tous
deux sur le même primitif sous-jacent (\texttt{scipy.spatial.cKDTree}),
afin que toute comparaison de coût computationnel isole la différence
\emph{algorithmique} (mise à l'échelle globale par paires contre
croissance locale de graphe) plutôt qu'une différence de qualité
d'implémentation entre une nouvelle méthode soigneusement optimisée et une
référence délibérément faible.

\subsection{Méthodologie de comparaison}

Pour un nuage de points et une dimension cible connue donnés, le nombre
minimal de points nécessaire à chaque méthode pour une convergence
\emph{stable} est déterminé~: le plus petit $n$ dans une grille testée pour
lequel l'estimation reste dans la tolérance de la cible \emph{et y reste}
pour tout $n$ supérieur testé --- se prémunissant contre un unique
franchissement fortuit d'une estimation encore mouvante, le mode d'échec
exact découvert dans les résultats initiaux sur les attracteurs chaotiques
(Section~\ref{sec:resultats}). Une « victoire » exige que les deux
méthodes convergent démontrablement~; une méthode qui ne converge jamais
ne peut se voir créditer d'avoir besoin de « moins » de points, ce qui
récompenserait l'échec de convergence comme s'il s'agissait d'efficacité.

\section{Banc d'Essai Initial : Dix Problèmes Physiques Connus}
\label{sec:resultats}

L'estimateur de dimension a été appliqué à dix problèmes couvrant des
cibles analytiques exactes, un théorème probabiliste rigoureux, et deux
valeurs de référence issues de la théorie du chaos classique.

\begin{table}[h]
\centering
\small
\begin{tabular}{@{}llll@{}}
\toprule
\# & Problème & Dimension connue & Source \\
\midrule
1 & Oscillateur harmonique & 1 (exact) & courbe fermée dans l'espace des phases \\
2 & Pendule non linéaire (grande amplitude) & 1 (exact) & orbite périodique fermée \\
3 & Orbite de Kepler à deux corps ($e=0,6$) & 1 (exact) & ellipse fermée \\
4 & Mars vs.\ données réelles JPL Horizons & 1 (exact) & arc lisse, éphéméride réelle \\
5 & Tore 2D quasi-périodique (Lissajous) & 2 (exact) & résultat standard de type KAM \\
6 & Mouvement brownien plan & 2 (exact) & théorème de Taylor (1953) \\
7 & Orbite périodique du problème restreint à 3 corps & 1 (exact, si fermeture vérifiée) & définition de cycle limite \\
8 & Attracteur de Lorenz & $\approx 2,05$ & Grassberger \& Procaccia (1983) \\
9 & Attracteur de R\"ossler & $\approx 2,01$ & Grassberger \& Procaccia (1983) \\
10 & Pendule amorti forcé (accroché en phase) & 1 (exact, si non chaotique) & définition de cycle limite \\
\bottomrule
\end{tabular}
\caption{Les dix problèmes du banc d'essai et leurs dimensions cibles issues de sources indépendantes.}
\end{table}

\subsection{Résultat affiché et pourquoi il est trompeur}

Les dix problèmes ont réussi leur tolérance déclarée. Un audit, suivi d'une
vérification indépendante de cet audit
(Section~\ref{sec:chaine-audit}), a établi que ce résultat affiché
surestime substantiellement ce qui a réellement été démontré~:

\begin{itemize}
\item \textbf{Six des dix « réussites » sont structurellement dégénérées.}
  Un graphe de $k$-PPV construit sur toute courbe fermée lisse est un
  réseau circulant en anneau exact, avec une taille de coquille
  parfaitement constante à chaque rayon. L'estimateur est
  mathématiquement incapable de retourner autre chose que la dimension 1
  avec $r^2=1,0$ sur une telle entrée --- l'ajustement parfait rapporté
  est une sentinelle signifiant « l'entrée n'avait aucune variation à
  ajuster », non une preuve que la méthode a mesuré quoi que ce soit avec
  précision. Cela affecte les problèmes 1, 2, 3, 4, 7 et 10.
\item \textbf{Aucun des deux résultats sur les attracteurs chaotiques
  n'était convergé en nombre de points.} Un balayage de convergence (en
  faisant varier le nombre de points échantillonnés, tout le reste étant
  fixé) a montré que les estimations de dimension de Lorenz et de
  R\"ossler augmentaient de façon monotone \emph{à travers} leurs valeurs
  de référence puis au-delà, continuant de croître au plus grand nombre
  de points testé. Les erreurs « réussies » rapportées (0,064 pour
  Lorenz, 0,014 pour R\"ossler) décrivaient où une quantité encore
  mouvante se trouvait échantillonnée, non une mesure convergée.
\item \textbf{Les deux résultats restants (tore de Lissajous, mouvement
  brownien)} étaient les moins dégénérés de l'ensemble, mais une
  correction ultérieure a montré qu'ils n'avaient pas été soumis à la
  même discipline de convergence/sensibilité à la grille
  d'échantillonnage appliquée aux cas chaotiques, et tous deux ont montré
  une sensibilité comparable au choix de la grille une fois vérifiés.
\end{itemize}

\subsection{Chaîne de relecture à plusieurs niveaux}
\label{sec:chaine-audit}

Dix agents indépendants ont chacun résolu un problème et rapporté un
résultat~; les dix auto-rapports se sont ensuite révélés numériquement
corrects. Un audit indépendant a examiné les dix résultats et a
correctement identifié les problèmes structurels ci-dessus. Une
\emph{seconde} relecture indépendante a alors été appliquée
spécifiquement à l'analyse écrite de l'audit lui-même --- et y a trouvé
trois erreurs réelles~: une généralisation excessive concernant la portée
d'un défaut logiciel (Section~\ref{sec:defauts}), une réattribution
spéculative ayant incorrectement infirmé une conclusion par ailleurs
correcte, et une affirmation factuellement inexacte quant aux
sous-résultats ayant reçu une vérification particulière, laquelle avait
étayé l'un des deux verdicts de « succès » phares de l'audit sans la même
rigueur appliquée ailleurs. Les trois erreurs sont documentées et
corrigées dans le registre du banc d'essai sous-jacent. Cette structure à
plusieurs niveaux --- exécutants, auditeur, et vérification indépendante
de l'auditeur lui-même --- constitue désormais la pratique permanente pour
tout résultat de ce programme destiné à être cité.

\section{Défauts Logiciels Identifiés et Corrigés}
\label{sec:defauts}

Deux défauts réels et indépendants dans l'implémentation même de
l'estimateur ont été trouvés comme conséquence directe de son application
à des problèmes réels plutôt qu'aux seuls réseaux synthétiques utilisés
dans ses tests unitaires d'origine.

\subsection{Annulation catastrophique du \texorpdfstring{$R^2$}{R2}}

Pour une séquence de coquilles parfaitement constante, la somme totale des
carrés dans le calcul du $R^2$ devrait être exactement nulle. Un
arrondi en virgule flottante dans le calcul intermédiaire de la moyenne la
laissait à environ $10^{-31}$ pour certaines longueurs d'ajustement et
valeurs de coquille au lieu d'exactement zéro, ce qui faisait prendre à la
comparaison à zéro exact de l'implémentation la branche du cas général et
diviser par ce résidu --- effondrant le $R^2$ rapporté à environ zéro par
pur bruit numérique, alors que l'estimation de dimension sous-jacente
restait correcte. Cela signifiait que les réglages par défaut documentés
de l'estimateur retournaient un résultat indéfini (\texttt{NaN}) pour
plusieurs réseaux ordinaires, y compris la configuration exacte utilisée
par quatre des dix problèmes du banc d'essai. La correction remplace la
comparaison à zéro exact par une tolérance relative à l'échelle~; vérifiée
sur chaque valeur de coquille trouvée déclenchant le défaut (6, 17, 18 à
la longueur d'ajustement concernée) ainsi que sur des contrôles négatifs
déjà corrects, avec deux tests de non-régression ajoutés.

\subsection{Défaut de performance dans la reconstruction d'adjacence}

Un défaut de performance distinct et sans rapport a été trouvé lors de la
construction du dispositif de comparaison~: la structure d'adjacence de
l'hypergraphe était reconstruite intégralement, à coût $O(\text{arêtes})$,
à chaque requête de distance, et l'estimateur de dimension l'interrogeait
une fois par rayon et par nœud échantillonné --- s'accumulant en milliers
de reconstructions redondantes pour un balayage de convergence réaliste.
Un balayage représentatif de 3200 points prenait 136 secondes avant
correction. La correction met en cache la structure d'adjacence par
instance d'hypergraphe (immuable) et restructure l'estimateur de dimension
pour effectuer une unique recherche en largeur incrémentale collectant les
volumes de tous les rayons en un seul passage, plutôt que de recalculer
depuis la source à chaque rayon. Le même balayage de 3200 points s'est
achevé en 5,0 secondes après correction --- une amélioration d'environ
27$\times$ --- avec des résultats numériquement \emph{identiques} avant et
après, confirmant que la correction n'a sacrifié aucune exactitude à la
vitesse. Cette amélioration de performance est la condition habilitante
spécifique de la nouvelle comparaison systématique rapportée comme en
cours à la Section~\ref{sec:etat-actuel}~: les balayages de convergence
qu'elle requiert n'étaient pas praticables avant cette correction.

\subsection{Détection des points dupliqués}

Une troisième amélioration, non pas une correction de bogue mais un ajout
de robustesse~: l'échantillonnage de plusieurs périodes d'une orbite
fermée produit des grappes de points quasi-identiques (séparés
d'environ $10^{-9}$ dans les cas observés), ce qui a silencieusement
corrompu trois des dix résultats initiaux du banc d'essai, avec une
sévérité allant d'une dégradation de la qualité d'ajustement jusqu'à un
échec complet de convergence. La construction $k$-PPV détecte désormais
les \emph{grappes} de doublons (et non de simples paires, puisqu'une
orbite à trois périodes produit des grappes de trois) via une analyse des
composantes connexes du graphe des paires en deçà du seuil, calibrée par
rapport à la diagonale de la boîte englobante du nuage de points plutôt
que par rapport à l'espacement local aux plus proches voisins --- une
première tentative utilisant ce dernier a échoué silencieusement, car le
plus proche voisin d'un point dupliqué \emph{est} son propre doublon,
corrompant la statistique même censée le détecter. Vérifiée sur la
reproduction exacte du défaut d'origine (300 points, chacun triplé par
période avec un décalage d'environ $10^{-9}$)~: l'implémentation corrigée
soulève une erreur explicite nommant les 300 grappes de taille trois, et
une option de nettoyage automatique récupère exactement les 300 points
d'origine et la dimension correcte de exactement 1,0.

\section{État Actuel de la Comparaison aux Méthodes Traditionnelles}
\label{sec:etat-actuel}

Une fois les défauts ci-dessus corrigés et une véritable référence de
méthode traditionnelle implémentée (Section 2.4), une nouvelle comparaison
systématique sur les dix problèmes du banc d'essai a été lancée, visant
l'objectif déclaré du programme~: une majorité (au moins huit sur dix) des
problèmes démontrant un avantage par rapport à la méthode de somme de
corrélation traditionnelle, défini par deux critères explicites et
vérifiables plutôt que par une notion vague unique de « meilleur »~:

\begin{enumerate}
\item \textbf{Économie de calcul}~: plus de 10\,\% de points
  d'échantillonnage en moins nécessaires pour une précision équivalente et
  \emph{stablement convergée}, contre la même tolérance.
\item \textbf{Robustesse à la densité} (la forme concrète et vérifiable de
  « l'évitement de singularité »)~: les trajectoires physiques présentant
  une approche rapprochée ou un point de rebroussement lent/rapide (un
  passage proche d'un primaire dans le problème restreint à trois corps,
  le passage au périhélie d'une orbite de Kepler excentrique, le point de
  rebroussement d'un pendule) sont échantillonnées bien plus densément
  dans le temps près de cette caractéristique qu'ailleurs, alors même que
  la courbe sous-jacente est uniformément unidimensionnelle partout. Une
  méthode \emph{globale} qui regroupe toutes les distances par paires en
  une seule statistique (la somme de corrélation traditionnelle) peut être
  mesurablement biaisée par une telle concentration de densité~; une
  méthode \emph{locale} qui n'examine jamais que les plus proches voisins
  de chaque point (l'estimateur par croissance de coquilles) ne l'est pas,
  par construction. Ceci est testé directement en comparant l'erreur de
  chaque méthode par rapport à la cible connue sur les données naturelles,
  non uniformément échantillonnées.
\end{enumerate}

\textbf{Cette comparaison était en cours au moment de la rédaction.} Un
point de données illustratif, vérifié manuellement, est disponible~: sur
un carré plein bidimensionnel (dimension cible 2, tolérance 0,15), la
méthode de somme de corrélation traditionnelle a atteint une convergence
stable à 100 points tandis que la méthode par croissance de coquilles en
a nécessité 400, un résultat favorisant la méthode traditionnelle sur ce
cas particulier ($-300\,\%$ en nombre de points relatif, c'est-à-dire
l'inverse d'une victoire poly-algébrique). Cette vérification ponctuelle
unique est rapportée par souci de transparence, non comme preuve du
résultat systématique sur l'ensemble des dix problèmes --- conformément à
la règle permanente de ce programme selon laquelle une affirmation reposant
sur une configuration unique, plutôt que sur une comparaison balayée et
convergée, n'est pas encore un résultat (Section 2.5). La comparaison
complète sur les dix problèmes, auditée et vérifiée indépendamment, sera
rapportée dans une révision ultérieure de ce document une fois achevée.

\section{Limitations}

\begin{itemize}
\item \textbf{L'estimateur n'a aucune liberté d'être inexact dans le
  régime dégénéré.} Toute courbe fermée lisse, échantillonnée densément
  et uniformément, sera lue comme de dimension 1 avec une qualité
  d'ajustement parfaite (mais non informative), indépendamment de la
  précision réelle de la méthode. Un indicateur \texttt{degenerate}
  distingue désormais explicitement ce cas dans le logiciel, mais tout
  résultat tiré des problèmes 1, 2, 3, 4, 7 ou 10 du banc d'essai initial
  doit être lu comme confirmant que le pipeline est correctement câblé,
  non comme une preuve de précision de mesure.
\item \textbf{Les cibles de dimension fractale (non entière) nécessitent
  un balayage de convergence, non une mesure unique}, et même le
  comportement de convergence du pipeline corrigé sur les deux
  attracteurs chaotiques de ce banc d'essai n'a pas encore été réétabli
  après la correction des défauts ci-dessus --- ceci fait partie du
  travail en cours de la Section~\ref{sec:etat-actuel}.
\item \textbf{La construction par $k$ plus proches voisins demeure une
  approximation}~: la distance de graphe approxime, sans lui être égale,
  la distance géodésique sur la variété sous-jacente, et cette
  approximation se dégrade près du bord d'une variété (documenté et fixé
  par un test de non-régression montrant une sous-estimation mesurable au
  coin d'un cube échantillonné) et est sensible au choix de $k$ et à la
  densité d'échantillonnage locale.
\item \textbf{Ce travail ne valide pas l'hypothèse physique de Wolfram.}
  Tout ce qui est rapporté ici est un énoncé sur le comportement d'un
  algorithme spécifique de théorie des graphes sur des nuages de points
  finis spécifiques --- Niveau B selon la convention épistémique de ce
  programme. Aucune affirmation n'est faite ni impliquée quant à savoir si
  la réécriture d'hypergraphes est une description fondamentale correcte
  de l'espace-temps physique~; cela demeure Niveau C (conjecture étiquetée,
  jamais structurante) tout au long de ce document.
\end{itemize}

\section{Discussion}

Le résultat le plus lourd de conséquences de cette étape est procédural
plutôt que purement technique~: un banc d'essai contre des cibles réelles
et indépendamment sourcées a trouvé de véritables défauts
(Section~\ref{sec:defauts}) que les tests unitaires propres de
l'estimateur, exercés uniquement contre des réseaux synthétiques conçus
pour être faciles, n'auraient pas pu trouver. Tout aussi lourd de
conséquences~: la \emph{relecture du banc d'essai} elle-même a nécessité
sa propre vérification indépendante avant que ses conclusions puissent
être considérées comme fiables --- un audit ayant correctement identifié
de véritables problèmes dans le travail initial n'était pas, sur cette
seule base, exempté de contenir lui-même de véritables problèmes. Ces
deux observations sont consignées comme leçons générales pour ce
programme (voir le document d'enseignements tirés associé) et sont
désormais intégrées au schéma de travail permanent utilisé pour tout
résultat de ce dépôt destiné à être cité~: mesure, audit, puis
vérification indépendante de l'audit, dans cet ordre, aucun des trois
n'étant considéré comme suffisant à lui seul.

\section{Conclusion et Prochaines Étapes}

La Phase 2, Étape 2 de ce programme a produit~: (i) une implémentation
fonctionnelle et testée de deux des quatre concepts de formalisation
envisagés pour la physique par hypergraphes~; (ii) un banc d'essai de dix
problèmes contre des cibles physiques indépendamment sourcées, dont le
résultat affiché a nécessité une correction substantielle sous audit mais
dont les défauts logiciels sous-jacents ont été véritablement trouvés et
véritablement corrigés~; (iii) une véritable référence classique et une
méthodologie de comparaison, remplaçant ce qui aurait autrement été une
affirmation de supériorité non falsifiable par deux critères spécifiques
et vérifiables. La comparaison systématique contre cette référence, sur
l'ensemble des dix problèmes, est la prochaine étape immédiate et était en
cours au moment de la rédaction. Son achèvement, audité et vérifié
indépendamment selon le même standard que tout autre résultat de ce
document, sera rapporté en addendum ou en révision plutôt que présumé à
l'avance.

\section*{Références}

\begin{itemize}
\item[] Wolfram, S. (2020). \emph{A Project to Find the Fundamental Theory of Physics.}
\item[] Grassberger, P. \& Procaccia, I. (1983). Measuring the strangeness of
  strange attractors. \emph{Physica D}, 9, 189--208.
\item[] Taylor, S. J. (1953). The Hausdorff $\alpha$-dimensional measure of
  Brownian paths in $n$-space. \emph{Proc. Camb. Phil. Soc.}, 49.
\item[] R\"ossler, O. (1976). An equation for continuous chaos.
  \emph{Phys. Lett. A}, 57.
\item[] Références internes~: \texttt{docs/HYPERGRAPH\_NOTES.md},
  \texttt{docs/POLY\_ALGEBRAIC\_BENCHMARK.md}, \texttt{docs/LL.md},
  \texttt{src/socrates/hypergraph/}.
\end{itemize}

\end{document}
